% Draft "Future work" subsection for NEUCOM-D-26-00713 (revision).
% Goes after \subsection{Limitations} in the Discussion.
% Two lengths: SHORT (~25%, active) and MEDIUM (~50%, commented out below).
% New citation keys are in latex/future_work_refs.bib (unused entries are harmless).

%% ---------- SHORT (~200 words) ----------
\subsection{Future work}
Three extensions follow directly from our results. First, our claims of
generalization rest on a narrow sample. Adding datasets from other modalities,
and sampling architecture, training set size and the weight target $w_{tgt}$
rather than fixing them, would show whether the effect of sleep holds for a
typical network as well as a typical dataset. Second, sleep, layer-wise
normalization and synaptic scaling reach similar accuracy but correct the weights
differently, so they may diverge under distribution shift, when rejecting
out-of-distribution inputs \citep{hendrycks_baseline_2017}, or when tasks are
learned in sequence \citep{tadros_sleep-like_2022}. Third, sleep should help most
when the learning rule has no check of its own on weight growth. Comparing it
across reward-modulated STDP \citep{fremaux_neuromodulated_2016}, e-prop
\citep{bellec_solution_2020} and surrogate gradient training
\citep{neftci_surrogate_2019} would show whether the benefit is specific to
Hebbian learning.

Two further questions concern the components of sleep. The noise we used was
unstructured, so its lack of effect does not rule out a role for structured
activity, such as replay of recently seen inputs \citep{deperrois_learning_2022}.
And because $n$ downscaling steps with factor $\lambda$ equal a single step with
factor $\lambda^{n}$, sleep could be implemented as one update per synapse, a
saving worth measuring on neuromorphic hardware \citep{davies_loihi_2018}.

%% ---------- MEDIUM (~430 words) ----------
% \subsection{Future work}
% Our results point to downscaling as the active ingredient of sleep, and to a
% benefit that holds across the datasets we sampled. The most direct extension is
% to widen that sample. Adding datasets from other modalities, such as natural
% images, spoken digits and event-based recordings, would sharpen the estimate of
% between-dataset variance. Architecture, training set size and the weight target
% $w_{tgt}$ could likewise be sampled rather than fixed, so that the reported
% effect describes a typical network as well as a typical dataset. Scaling up also
% makes the cost of sleep relevant: because $n$ downscaling steps with factor
% $\lambda$ equal a single step with factor $\lambda^{n}$
% (Appendix~\ref{sec:app:mech_weight_decay}), sleep could be implemented as one
% update per synapse, a saving worth measuring on neuromorphic hardware
% \citep{davies_loihi_2018}.
%
% We evaluated all models on test images from the training distribution. Sleep,
% layer-wise normalization and synaptic scaling reached similar accuracy, but
% normalization rescales the weights to a fixed total, whereas sleep compresses
% their distribution toward the target. These differences may matter under
% distribution shift, or when the network must reject inputs from an unseen
% dataset, for which the spike count of the most responsive class-assigned
% neurons offers a natural signal \citep{hendrycks_baseline_2017}. Learning tasks
% in sequence is a related test \citep{tadros_sleep-like_2022}: pulling all
% weights toward a common target may protect earlier representations, or erase
% them faster.
%
% The ablation leaves the role of noise open. The noise we used was unstructured
% and may carry no information for STDP to strengthen; replacing it with
% structured activity, such as replay of recently seen inputs
% \citep{deperrois_learning_2022}, would test whether sleep-phase plasticity
% contributes once it has something to consolidate.
%
% Finally, sleep should help most when the learning rule has no check of its own
% on weight growth, as with unsupervised STDP. Surrogate gradient training
% \citep{neftci_surrogate_2019} corrects weights whenever their growth harms the
% objective and may gain little from sleep, while reward-modulated STDP
% \citep{fremaux_neuromodulated_2016} and e-prop \citep{bellec_solution_2020} lie
% between the two. Comparing sleep with normalization across these rules would
% show whether its benefit is specific to Hebbian learning.
